\documentclass[11pt]{amsart}
\usepackage[T1]{fontenc}
\usepackage{amsmath,amssymb,amsthm}
\usepackage{graphicx}
\usepackage[colorlinks=true,linkcolor=blue,citecolor=blue,urlcolor=blue]{hyperref}

\newtheorem{theorem}{Theorem}
\newtheorem{lemma}{Lemma}
\newtheorem{remark}{Remark}

\newcommand{\dist}{\operatorname{dist}}
\newcommand{\R}{\mathbb R}
\newcommand{\elltwo}{\ell_2}

\begin{document}

\title{A Non-Uniformly Smooth Space With Property \(\mathcal D(s,c_1)\)}
\author{Run fa\_banach\_001}
\date{2026-06-16}
\maketitle

\section*{Status}

This is a candidate full answer to Remark 2.11 of \cite{DGHKT}. The answer is
positive: property \(\mathcal D(s,c_1)\) can hold for a non-uniformly smooth
Banach space and a dictionary.

\section*{Source Question}

Definition 1.12 of \cite{DGHKT} says that \((X,\|\cdot\|,\mathcal D)\) has
property \(\mathcal D(s,c_1)\) if, for some \(s>1\) and \(c_1>0\),
\[
  \dist(f,[\varphi])
  \leq \|f\|\bigl(1-c_1 |F_f(\varphi)|^s\bigr)
\]
for every \(f\neq 0\) and every \(\varphi\in\mathcal D\), where \(F_f\) is an
associated norming functional. Remark 2.11 asks whether this property could
hold for some non-uniformly smooth Banach space and some dictionary.

\begin{center}
\includegraphics[width=\textwidth]{figures/open_problem_crop.png}
\end{center}

\section*{Proof Intuition}

The example is a Hilbert norm with a small absolute-value kink:
\[
  N_\varepsilon(x,y)=(x^2+y^2)^{1/2}+\varepsilon |x+y|.
\]
The kink makes the norm non-smooth, hence not uniformly smooth. But the kink is
transverse to the two coordinate dictionary lines. For each coordinate line,
the best distance in the \(N_\varepsilon\)-norm is obtained by a one-variable
convex minimization. The distance defect from this minimum is quadratic in the
chosen norming functional's value on the coordinate atom. This gives
\(\mathcal D(2,c)\) on the two-dimensional block. Taking an \(\ell_2\)-sum
with a Hilbert space preserves this coordinate-dictionary estimate and gives
an infinite-dimensional example.

\section*{The Example}

Fix \(0<\varepsilon<1/2\). On \(\R^2\), define
\[
  N_\varepsilon(x,y)=\sqrt{x^2+y^2}+\varepsilon |x+y|.
\]
Let \(E_\varepsilon=(\R^2,N_\varepsilon)\), and let \(H=\elltwo\). Set
\[
  X=E_\varepsilon\oplus_2 H,\qquad
  \|(z,h)\|_X=\bigl(N_\varepsilon(z)^2+\|h\|_2^2\bigr)^{1/2}.
\]
Let
\[
  u_1=\frac{(1,0)}{1+\varepsilon},\qquad
  u_2=\frac{(0,1)}{1+\varepsilon},
\]
which are unit vectors in \(E_\varepsilon\), and let \((e_n)\) be the usual
unit vector basis of \(H\). The dictionary is
\[
  \mathcal D=\{(u_1,0),(u_2,0)\}\cup\{(0,e_n):n\geq1\}.
\]
Its linear span is dense in \(X\).

\section*{Finite-Dimensional Block Estimate}

For \(z=(x,y)\neq0\), put \(r=(x^2+y^2)^{1/2}\). Choose the following norming
functional for \(E_\varepsilon\):
\[
  G_z(a,b)=\frac{xa+yb}{r}+\varepsilon\theta_z(a+b),
\]
where \(\theta_z=\operatorname{sign}(x+y)\) if \(x+y\neq0\), and
\(\theta_z=0\) if \(x+y=0\). This is a subgradient of \(N_\varepsilon\) at
\(z\), hence \(G_z(z)=N_\varepsilon(z)\) and \(\|G_z\|_{E_\varepsilon^*}=1\).

\begin{lemma}\label{lem:block}
There is a constant \(c_E=c_E(\varepsilon)>0\) such that, for \(j=1,2\),
\[
  \dist_{E_\varepsilon}(z,[u_j])
  \leq N_\varepsilon(z)\bigl(1-c_E|G_z(u_j)|^2\bigr),
  \qquad z\neq0.
\]
\end{lemma}

\begin{proof}
We give the proof for \(j=1\); the other case is symmetric. Since
\([u_1]=[(1,0)]\), the distance from \(z=(x,y)\) to this line is
\[
  \delta_1(z)=\inf_{\lambda\in\R} N_\varepsilon(x-\lambda,y)
  =\inf_{t\in\R}\left(\sqrt{t^2+y^2}+\varepsilon |t+y|\right).
\]
If \(y=0\), then \(\delta_1(z)=0\) and the desired inequality is immediate
after decreasing \(c_E\leq1\). Assume \(y\neq0\), write
\(\eta=\operatorname{sign} y\), and set \(t=x/|y|\). By homogeneity,
\[
  \delta_1(z)=|y|m_\eta,\qquad
  N_\varepsilon(z)=|y|h_\eta(t),
\]
where
\[
  h_\eta(t)=\sqrt{t^2+1}+\varepsilon |t+\eta|,\qquad
  m_\eta=\inf_{v\in\R} h_\eta(v).
\]
For \(0<\varepsilon<1/2\), the unique minimizer of \(h_\eta\) is
\[
  t_\eta=-\eta\,\frac{\varepsilon}{\sqrt{1-\varepsilon^2}},
\]
which is not the kink point \(-\eta\). Away from the kink, the selected value
\((1+\varepsilon)G_z(u_1)\) is exactly the derivative \(h_\eta'(t)\). At the
kink our selected value is one element of the subdifferential, while
\(h_\eta(-\eta)>m_\eta\).

Consider, where the denominator is nonzero,
\[
  Q_\eta(t)=
  \frac{h_\eta(t)-m_\eta}
       {h_\eta(t)\left|h_\eta'(t)/(1+\varepsilon)\right|^2},
\]
using the selected subgradient value at the kink. This quotient has a positive
lower bound. Indeed, outside a compact interval it has positive liminf; away
from \(t_\eta\) and the kink this is ordinary continuity; at the kink the
numerator is strictly positive; and near \(t_\eta\),
\[
  h_\eta(t)-m_\eta
    =\tfrac12 h_\eta''(t_\eta)(t-t_\eta)^2+o((t-t_\eta)^2),\qquad
  h_\eta'(t)=h_\eta''(t_\eta)(t-t_\eta)+o(t-t_\eta),
\]
with \(h_\eta''(t_\eta)>0\). Hence \(Q_\eta\) extends there with a positive
finite limit. Taking the minimum over \(\eta=\pm1\) gives \(c_E>0\). Rewriting
the inequality \(Q_\eta(t)\geq c_E\) in homogeneous variables proves the
claim for \(u_1\).
\end{proof}

\section*{Hilbert Sum Stability}

The Hilbert basis has the corresponding estimate with exponent \(2\): if
\(a_n=\langle h,e_n\rangle/\|h\|_2\), then
\[
  \dist_H(h,[e_n])
  =\|h\|_2(1-|a_n|^2)^{1/2}
  \leq \|h\|_2(1-\tfrac12 |a_n|^2).
\]

Choose norming functionals on \(X\) as follows. If \(z\neq0\) and \(h\neq0\),
with \(R=\|(z,h)\|_X\), set
\[
  F_{(z,h)}(w,k)=
  \frac{N_\varepsilon(z)}{R}G_z(w)
  +\frac{\|h\|_2}{R}\left\langle \frac{h}{\|h\|_2},k\right\rangle.
\]
The evident one-component formula is used when \(z=0\) or \(h=0\).

\begin{theorem}
The Banach space \(X=E_\varepsilon\oplus_2\elltwo\) is not uniformly smooth,
but \((X,\mathcal D)\) has property \(\mathcal D(2,c)\) for some
\(c=c(\varepsilon)>0\).
\end{theorem}

\begin{proof}
The norm \(N_\varepsilon\) is not G\^ateaux differentiable at any nonzero
point of the line \(x+y=0\), because the term \(|x+y|\) has a kink there.
Therefore the norm of \(X\) is not G\^ateaux differentiable at
\(((1,-1),0)\). Uniform smoothness implies G\^ateaux differentiability away
from zero, so \(X\) is not uniformly smooth.

It remains to prove property \(\mathcal D(2,c)\). Let \(R=\|(z,h)\|_X\),
\(A=N_\varepsilon(z)\), and \(B=\|h\|_2\). First take an atom
\(\varphi=(u_j,0)\). Put \(a=G_z(u_j)\), with \(a=0\) if \(z=0\). By
Lemma \ref{lem:block},
\[
  \dist_{E_\varepsilon}(z,[u_j])\leq A(1-c_E|a|^2).
\]
Hence
\[
\begin{aligned}
  \dist_X((z,h),[\varphi])^2
  &\leq A^2(1-c_E|a|^2)^2+B^2  \\
  &\leq R^2\left(1-c_E\frac{A^2}{R^2}|a|^2\right).
\end{aligned}
\]
Since \(|F_{(z,h)}(\varphi)|=(A/R)|a|\), taking square roots yields
\[
  \dist_X((z,h),[\varphi])
  \leq R\left(1-\frac{c_E}{2}|F_{(z,h)}(\varphi)|^2\right).
\]

For a Hilbert atom \(\varphi=(0,e_n)\), the same computation, using the
Hilbert estimate above in place of Lemma \ref{lem:block}, gives
\[
  \dist_X((z,h),[\varphi])
  \leq R\left(1-\frac14 |F_{(z,h)}(\varphi)|^2\right).
\]
Thus property \(\mathcal D(2,c)\) holds with
\[
  c=\min\{c_E/2,1/4\}>0.
\]
\end{proof}

\section*{History: Why the First Obstruction Was Not the End}

Before finding the positive example above, the direct attack produced a useful
negative result for a natural class of nonsmooth spaces. It is included here
so the full packet records the route cleanly.

\begin{theorem}[Historical obstruction]
Let \(E\) and \(F\) be nonzero Banach spaces, and let
\[
  Y=E\oplus_1 F,\qquad \|(e,f)\|=\|e\|+\|f\|.
\]
Then no nonempty dictionary in \(Y\) can satisfy property
\(\mathcal D(s,c)\) for any \(s>1\) and \(c>0\).
\end{theorem}

\begin{proof}
It is enough to show that every unit atom \(\varphi=(u,v)\in Y\) fails the
one-atom inequality. Put \(A=\|u\|\), \(B=\|v\|\), so \(A+B=1\).

Assume first \(A,B>0\), and write \(e=u/A\), \(y=v/B\). If \(A\leq B\), take
\[
  x_t=(e,t y),\qquad t>0.
\]
Then \(\|x_t\|=1+t\). Every norming functional for \(x_t\) evaluates
\(\varphi\) as \(1\). On the other hand, for any scalar \(\lambda\),
\[
  \|x_t-\lambda\varphi\|
  = |1-\lambda A|+|t-\lambda B|
  \geq |1-\lambda A|+\frac AB |t-\lambda B|
  \geq 1-\frac AB t,
\]
and equality is attained at \(\lambda=t/B\). Hence
\[
  \frac{\dist(x_t,[\varphi])}{\|x_t\|}
  =\frac{1-(A/B)t}{1+t}\longrightarrow 1,
  \qquad |F_{x_t}(\varphi)|=1.
\]
This contradicts any estimate
\(\dist(x,[\varphi])\leq \|x\|(1-c|F_x(\varphi)|^s)\) with \(c>0\).
If \(A>B\), the same proof uses \(x_t=(t e,y)\).

If \(A=0\), choose any unit \(e\in E\) and take \(x_t=(e,t v)\). Then
\(\dist(x_t,[\varphi])=1\), \(\|x_t\|=1+t\), and every norming functional
again evaluates \(\varphi\) as \(1\). Letting \(t\downarrow0\) gives the same
contradiction. The case \(B=0\) is symmetric.
\end{proof}

This obstruction explains why the successful example uses a small Hilbertian
kink rather than an \(\ell_1\)-sum corner. The full packet's history folder
contains the original partial packet and its superseded ledger record.

\section*{Verification And Novelty Notes}

The proof above is independent of numerical experimentation. As a sanity
check, the script \texttt{code/sample\_block\_ratio.py} samples the quotient
in Lemma \ref{lem:block} for \(\varepsilon=1/4\); it found a positive sampled
minimum approximately \(0.23638\).

Local run indexes were searched for arXiv:1909.13536, the exact title,
\(\mathcal D(s,c_1)\), non-uniformly smooth spaces, and related WCGA terms.
Web searches for the exact Remark 2.11 phrasing and property
\(\mathcal D(s,c_1)\) found the source arXiv page \cite{DGHKT} but no later
paper explicitly answering the question. Novelty confidence is therefore
reasonable but still subject to human literature review. The previous
\(\ell_1\)-sum obstruction is no longer indexed as a separate partial result;
it has been packed into this full result as historical context.

\begin{thebibliography}{9}

\bibitem{DGHKT}
S. Dilworth, G. Garrig\'os, E. Hern\'andez, D. Kutzarova, and V. Temlyakov,
\emph{Lebesgue-type inequalities in greedy approximation},
arXiv:1909.13536.

\end{thebibliography}

\end{document}
